%قبل از شروع مسئله ترجیحاً چیزی ننویسید. اما اگر خواستید بنویسید، در قالب لیستی (فقرات و شمارش) بنویسید.
%برای اضافه کردن نمره، بعدی از شروع‌مسئله، [x] را اضافه کنید و به جای x  نمره مورد نظر را قرار دهید. برای زیرمسئله‌ها هم به همین صورت است.
%\شروع‌مسئله[10]
\شروع‌مسئله
% صورت اصلی مسئله را در این‌جا بنویسید.
در این سوال می‌خواهیم داده‌ساختاری شبیه هرم
\LTRfootnote{heap}
را پیاده‌سازی و تحلیل کنیم. با این تفاوت که برخلاف Heap که هر راس حداکثر دو فرزند داشت، در این جا هر راس می‌توانند حداکثر $d$ فرزند داشته باشد. به این داده ساختار هرم \lr{d}-تایی
می‌گوییم.
%اگر مسئله چند قسمت داشت به این صورت عمل کنید:
\شروع{زیرمسائل}

\زیرمسئله
چنین هرمی که هر راس‌‌آن حداکثر $d$ فرزند دارد را چه طور در آرایه ذخیره می‌کنید؟ برای حل این سوال کافیست توابعی که اندیس فرزند $k$ ام عنصر $i$ ام و والد عنصر $i$ ام را برمی‌گردانند را بنویسید.

\زیرمسئله
ارتفاع یک هرم \lr{d}-تایی با $n$ عضو چند است؟

\زیرمسئله
یک پیاده‌سازی بهینه برای تابع $Extract\_Max$ در هرم بیشینه \lr{d}-تایی ارائه کنید و آن را برحسب $n$ و $d$ از لحاظ زمانی تحلیل نمایید.

\زیرمسئله
یک پیاده‌سازی بهینه برای تابع $Insert$ در هرم بیشینه \lr{d}-تایی ارائه کنید و آن را برحسب $n$ و $d$ از لحاظ زمانی تحلیل نمایید.

\زیرمسئله
یک پیاده‌سازی بهینه برای تابع $Increase\_Key(A,i,k)$ در هرم بیشینه \lr{d}-تایی ارائه کنید. این تابع در صورتی که عنصر $i$ ام بزرگ‌تر از $k$ باشد خطا بر می‌گرداند و در غیر این صورت مقدار $A[i]$ را برابر $k$ قرار داده و سپس هرم را آپدیت می‌کند تا خاصیت هرم بیشینه بودن آن حفظ شود. الگوریتم را برحسب $n$ و $d$ از لحاظ زمانی تحلیل نمایید.



\پایان{زیرمسائل}
\solution{

برای قسمت اول توابع را به این شکل پیاده سازی می‌کنیم.

\begin{latin}
	\begin{algorithm}[H]
		\begin{algorithmic}
			
			\Procedure {$Parent(i)$}{}
			\Return $\lfloor \frac{i}{d} \rfloor $
			\EndProcedure
			
			\Procedure {$Child(i,k)$}{}
			 \Return $d(i-1) +1 + k$
			\EndProcedure
			
			
		\end{algorithmic}
	\end{algorithm}	
\end{latin} 




برای قسمت دوم جواب
$\Theta(\log_{d}(n))$
است.


برای قسمت سوم دو تابع 
$HeapExtractMin$
و
$DMaxHeapify$
را پیاده‌سازی می‌کنیم.

\begin{latin}
	\begin{algorithm}[H]
		\begin{algorithmic}
			
			\Procedure {$HeapExtractMin(A)$}{}
			\State $max = A[1]$
			\State $A[1] = A[A.HeapSize]$
			\State $A.HeapSize = A.HeapSize-1$
			\State $DMaxHeapify(A,1)$
			\EndProcedure
			
			\Procedure {$DMaxHeapify(A,i)$}{}
			\State $largest = i$
			\For{$k=1$ to $d$}
			\If {$Child(i,k) \leq A.HeapSize$ and $A[Child(i,k)] > A[i]$}
			\If {$A[Child(i,k)] > largest$}
			\State $largest = A[Child(i,k)]$ 
			\EndIf
			\EndIf
			\EndFor
			\If {$largest \neq i$}
			\State swap $A[i]$ with $A[largest]$
			\State $DMaxHeapify(A,largest)$
			\EndIf
			\EndProcedure
			
			
		\end{algorithmic}
	\end{algorithm}	
\end{latin} 

برای تحلیل زمانی، تابع $HeapExtractMin$ تعدادی عملیات با مرتبه زمانی ثابت و یک بار فراخوانی $DMaxHeapify$ است. برای این تابع، تعداد دفعات فراخوانی به عمق هرم بستگی دارد و در هر بار اجرا، یک حلقه $d$ تایی داریم پس پیچیدگی زمانی آن از 
$O(d \log_{d}(n))$
است.


برای قسمت بعدی تابع $Insert$ را پیاده‌سازی می‌کنیم.


\begin{latin}
	\begin{algorithm}[H]
		\begin{algorithmic}
			
			\Procedure {$Insert(A,key)$}{}
			\State $A.HeapSize = A.HeapSize + 1$
			\State $A[A.HeapSize] = key$
			\State $i = A.HeapSize$
			\While {$i>1$ and $A[Parent(i)] < A[i]$}
			\State swap $A[i]$ with $A[Parent(i)]$
			\State $i = Parent(i)$
			\EndWhile
			\EndProcedure
			
			
		\end{algorithmic}
	\end{algorithm}	
\end{latin} 

مرتبه زمانی اجرای این الگوریتم وابسته به عمق هرم است یعنی این الگوریتم از
$O(\log_{d}(n))$
است.

برای قسمت آخر تابع $IncreaseKey$ را به صورت زیر پیاده سازی می‌کنیم:

\begin{latin}
	\begin{algorithm}[H]
		\begin{algorithmic}
			
			\Procedure {$IncraseKey(A,i,k)$}{}
			\If {$k < A[i]$}
			\State Throw Error
			\EndIf
			\State $A[i] = k$
			\While {$i>1$ and $A[Parent(i)] < A[i]$}
			\State swap $A[i]$ with $A[Parent(i)]$
			\State $i = Parent(i)$
			\EndWhile
			\EndProcedure
			
			
		\end{algorithmic}
	\end{algorithm}	
\end{latin} 

مشابه بخش قبل مرتبه زمانی اجرای این الگوریتم وابسته به عمق هرم است یعنی این الگوریتم از
$O(\log_{d}(n))$
است.

}



%%%%%%%%%%%%%%%%%%%%%%%%%%%%%%%%%%%%%%%%%%%%%%%%%%%%%
% برخی دستورات مفید برای تمرین‌ها. %
%%%%%%%%%%%%%%%%%%%%%%%%%%%%%%%%%%%%%%%%%%%%%%%%%%%%%

% فرمول‌های نمونه:
%$f(x) = \Theta(x)$
%$T(x, y) \leq \log(x) + 2^y + \binom{2n}{n}$
% {\tt code\_function}


% لیست موارد بدون شماره‌گذاری:
%\شروع{فقرات}
%    \فقره مورد اول 
%        \شروع{فقرات}
%            \فقره مورد اول داخلی 
%                \شروع{فقرات}
%                    \فقره اولین مورد خیلی داخلی 
%                    \فقره دومین مورد خیلی داخلی 
%                \پایان{فقرات}
%            \فقره مورد دوم داخلی 
%        \پایان{فقرات}
%    \فقره مورد دوم 
%\پایان{فقرات}

% لیست موارد شماره‌گذاری شده:
%\شروع{شمارش}
%    \فقره یک مورد 
%    \فقره یک مورد دیگر 
%    \فقره باز هم یک مورد دیگر
%\پایان{شمارش}

% برای نوشتن یک فرمول چند خطی که یک شماره بخوره. برای حالت‌های دیگر به سایت زیر مراجعه کنید:
% https://www.overleaf.com/learn/latex/Aligning_equations_with_amsmath
%\begin{align} 
%    \begin{split}
%        (x+y)^3 &= (x+y)^2(x+y) \\
%                &= (x^2+2xy+y^2)(x+y) \\
%                &= (x^3+2x^2y+xy^2) + (x^2y+2xy^2+y^3) \\
%                &= x^3+3x^2y+3xy^2+y^3
%    \end{split}                                 
%\end{align}

% نوشتن ماتریس و آرایه دوبعدی
%\begin{align}
%    A = 
%    \begin{bmatrix}
%        A_{11} & A_{21} \\
%        A_{21} & A_{22}
%    \end{bmatrix}
%\end{align}

% اضافه کردن تصویر با شرح. برای اضافه کردن تصویر بدون شرح، فقط دستور درج‌تصویر را استفاده کنید.
%\شروع{شکل}[h!]
%\تنظیم‌ازوسط
%  \درج‌تصویر[پهنا=0.5\پهنای‌متن]{img.jpg}
%\شرح{شرح تصویر}
%\برچسب{تصویر:۱}
%\پایان{شکل}
%\پایان{زیرمسائل}

%کد زیر را ببینید.
%%\lstinputlisting{source_filename.py}
%\begin{latin}
%\begin{lstlisting}
%# Import the modules
%import sys
%import random
%
%ans = True
%
%while ans:
%question = raw_input("Ask the magic 8 ball a question: (press enter to quit) ")
%
%answers = random.randint(1,8)
%
%if question == "":
%sys.exit()
%
%elif answers == 1:
%print "It is certain"
%
%elif answers == 2:
%print "Outlook good"
%
%elif answers == 3:
%print "You may rely on it"
%
%elif answers == 4:
%print "Ask again later"
%
%elif answers == 5:
%print "Concentrate and ask again"
%
%elif answers == 6:
%print "Reply hazy, try again"
%
%elif answers == 7:
%print "My reply is no"
%
%elif answers == 8:
%print "My sources say no"
%\end{lstlisting}
%\end{latin}
